\section*{\S\ 4. Die Zerlegungsidentit\"aten.}
Wir betrachten nun den allgemeinsten Fall der Identit\"aten zwischen beliebigen Symbol- und Variabelnreihen, in deren Schlu\ss ausdr\"ucken ohne Vornahme der Substitution (11.) eine Reihe $p_\tau$ in ihre einzelnen Reihen $y^{(1)}\cdots y^{(\tau)}$ zerlegt auftritt, und bezeichnen diese Identit\"aten als ,,Zerlegungsidentit\"aten.`` Im Anschlu\ss\ an die Bemerkung am Schlu\ss\ des vorigen Paragraphen bestimmen wir sie vorerst f\"ur den Fall nur zweier Symbolreihen, entwickeln also den Ausdruck $\pair{q_\rho A^\sigma}{B_{\rho+\sigma-\tau}p_\tau}$. Wir setzen dabei:
\begin{equation}
\begin{aligned}
q_{\rho-\alpha}^{(i_1\ldots i_\alpha)}&\sim p_{n-\rho}\,y^{(i_1)}\cdots y^{(i_\alpha)},\qquad
p_{\tau-\alpha}^{(i_1\ldots i_\alpha)}=y^{(i_{\alpha+1})}\cdots y^{(i_\tau)},\qquad
p_\tau=y^{(1)}\cdots y^{(\tau)}.
\end{aligned}
\tag{22}
\end{equation}
In bezug auf die Zahlen $\rho,\sigma,\tau$ sind vier F\"alle zu unterscheiden:
\[
\begin{array}{llll}
1)&\rho+\sigma\leqq n,& \tau\leqq\rho,& \tau\leqq\sigma;\\[2pt]
2)&\rho+\sigma\leqq n,& \tau\geqq\rho,& \tau\leqq\sigma;\\[2pt]
3)&\rho+\sigma\leqq n,& \tau\leqq\rho,& \tau>\sigma;\\[2pt]
4)&\rho+\sigma\leqq n,& \tau\geqq\rho,& \tau\geqq\sigma,
\end{array}
\]
von denen wir den ersten herausgreifen wollen, um danach die \"ubrigen kurz zu charakterisieren. Nach (20.), (10.) und (9.) kommt, wenn wir die Reihe $y^{(1)}y^{(2)}\cdots y^{(\tau)}B_{\rho+\sigma-\tau}$ in $y^{(1)}$ und $y^{(2)}\cdots y^{(\tau)}B_{\rho+\sigma-\tau}$ zerlegen:
\begin{equation}
\begin{aligned}
\pair{q_\rho A^\sigma}{y^{(1)}y^{(2)}\cdots y^{(\tau)}B_{\rho+\sigma-\tau}}
&=\pair{q_{\rho-1}^{(1)}A^\sigma}{y^{(2)}\cdots y^{(\tau)}B_{\rho+\sigma-\tau}}\\
&\quad+(-1)^\rho\pair{q_\rho[A_{n-\sigma}y^{(1)}]^{\sigma-1}}{y^{(2)}\cdots y^{(\tau)}B_{\rho+\sigma-\tau}}.
\end{aligned}
\tag{23}
\end{equation}
Das letzte Glied in (23.) hat nun genau die Form des urspr\"unglichen Ausdrucks -- es ist nur $\sigma$ durch $\sigma-1$, $\tau$ durch $\tau-1$ ersetzt --, l\"a\ss t also wieder die Gleichung (23.) zu. Wir k\"onnen also, da $\tau\leqq\sigma$, die Operation (23.) $\tau$-mal anwenden und gelangen unter Ber\"ucksichtigung von (10.) zu der Relation:
\begin{equation}
\begin{aligned}
\pair{q_\rho A^\sigma}{p_\tau B_{\rho+\sigma-\tau}}
&=(-1)^{\rho\sigma+(\sigma+\tau)(n+1)}
  \pair{q_\rho B^{n-\rho-\sigma+\tau}}{A_{n-\sigma}p_\tau}\\
&\quad+\sum_{i_1=1}^{\tau}(-1)^{\rho(i_1-1)}
 \pair{q_{\rho-1}^{(i_1)}[A_{n-\sigma}y^{(1)}\cdots y^{(i_1-1)}]^{\sigma-i_1+1}}
 {y^{(i_1+1)}\cdots y^{(\tau)}B_{\rho+\sigma-\tau}}.
\end{aligned}
\tag{24}
\end{equation}
Ein jedes Glied unter dem Summenzeichen ist nun wieder vom Typus des urspr\"unglichen Ausdrucks; es ist $\rho$ durch $\rho-1$, $\sigma$ durch $\sigma-i_1+1$, $\tau$ durch $\tau-i_1$ ersetzt, die Bedingung 1) also noch st\"arker erf\"ullt. Wir k\"onnen deshalb, da $\tau\leqq\rho$, die Operation (24.) $\tau$-mal anwenden und erhalten die gew\"unschte Zerlegungsidentit\"at:
\begin{equation}
\begin{aligned}
\pair{q_\rho A^\sigma}{p_\tau B_{\rho+\sigma-\tau}}
&=\eps_{i_0}\pair{q_\rho B^{n-\rho-\sigma+\tau}}{A_{n-\sigma}p_\tau}\\
&\quad+\sum_{i_1=1}^{\tau}\eps_{i_1}
 \pair{q_{\rho-1}^{(i_1)}B^{n-\rho-\sigma+\tau}}{A_{n-\sigma}p_{\tau-1}^{(i_1)}}\\
&\quad+\sum_{i_2=i_1+1}^{\tau}\eps_{i_2}
 \pair{q_{\rho-2}^{(i_1i_2)}B^{n-\rho-\sigma+\tau}}{A_{n-\sigma}p_{\tau-2}^{(i_1i_2)}}+\cdots\\
&\quad+\eps_{i_\tau}\pair{q_{\rho-\tau}^{(i_1\ldots i_\tau)}B^{n-\rho-\sigma+\tau}}{A_{n-\sigma}}\\
&=\sum_{\alpha=0}^{\tau}\sum_{i_\alpha=i_{\alpha-1}+1}^{\tau}\eps_{i_\alpha}
 \pair{q_{\rho-\alpha}^{(i_1\ldots i_\alpha)}B^{n-\rho-\sigma+\tau}}{A_{n-\sigma}p_{\tau-\alpha}^{(i_1\ldots i_\alpha)}}.
\end{aligned}
\tag{25}
\end{equation}
F\"ur den Modul $\eps$ ergibt sich aus (24.):
\begin{equation}
\begin{aligned}
\eps_{i_\alpha}&=(-1)^{\delta_{i_\alpha}},\qquad
\delta_{i_\alpha}=\sum_{\beta=1}^{\alpha}(\rho-\beta+1)(i_{\beta-1}+i_\beta+1)\\
&\quad +(\rho-\alpha)(\sigma+i_\alpha+\alpha)+(\sigma+\tau+\alpha)(n+1),
\end{aligned}
\tag{26}
\end{equation}
wo $i_0=0$ zu setzen ist. Das Summenzeichen in (25.) ist so zu verstehen, da\ss\ jeder Kombination $i_1i_2\ldots i_{\alpha-1}$, wo $i_1<i_2<\cdots<i_{\alpha-1}$, alle Werte $i_\alpha$, f\"ur die $i_{\alpha-1}<i_\alpha\leqq\tau$, zugef\"ugt werden. Die einzelnen Summen in (25.) gen\"ugen, wie man sich unter Ber\"ucksichtigung von (26.) leicht \"uberzeugt, den (19.) analogen Differentialgleichungen:
\begin{equation}
 \left(\frac{\partial}{\partial y^{(i)}}\middle|y^{(i+1)}\right)=0,
 \qquad (i=1,2,\ldots,\tau-1),
\tag{27}
\end{equation}
h\"angen also nur von der Reihe $p_\tau$ ab. Die Zerlegungsidentit\"at ist also eine unzerlegbare Identit\"at zwischen den Reihen $A,B,q,p$; auf die explizite Darstellung kommen wir im n\"achsten Paragraphen zur\"uck.

Im Fall 2) l\"a\ss t sich die Operation (23.) ebenfalls $\tau$-mal anwenden, die Operation (24.) aber bricht nach dem $\rho$-ten Male ab. Es ist also im Schlu\ss ausdruck von (25.) das erste Summenzeichen zu ersetzen durch $\sum_{\alpha=0}^{\rho}$. Im Fall 3) und 4) l\"a\ss t sich die Operation (23.) nur $\sigma$-mal ausf\"uhren; an Stelle von (24.) tritt dann die Relation:
\begin{equation}
\begin{aligned}
\pair{q_\rho A^\sigma}{p_\tau B_{\rho+\sigma-\tau}}
&=(-1)^{\rho\sigma}
 \pair{q_{\rho-(\tau-\sigma)}^{(\sigma+1\ldots\tau)}B^{n-\rho+\tau-\sigma}}{A_{n-\sigma}p_{\tau-(\tau-\sigma)}^{(\sigma+1\ldots\tau)}}\\
&\quad+\sum_{i_1=1}^{\sigma}(-1)^{\rho(i_1-1)}
 \pair{q_{\rho-1}^{(i_1)}[A_{n-\sigma}y^{(1)}\cdots y^{(i_1-1)}]^{\sigma-i_1+1}}
 {y^{(i_1+1)}\cdots y^{(\tau)}B_{\rho+\sigma-\tau}};
\end{aligned}
\tag{28}
\end{equation}
die $(\tau-\sigma)$-malige Wiederholung von (28.), wobei das Summenzeichen in $\sum_{i_\beta=i_{\beta-1}+1}^{\sigma+\beta-1}$ \"ubergeht -- wie die $(\sigma-i_1+1)$-malige Anwendung von (23.) auf die Glieder unter dem Summenzeichen in (28.) zeigt --, gibt alle Glieder der dem Index $\alpha=\tau-\sigma$ entsprechenden Einzelsumme von (25.) mit entsprechendem Vorzeichen. Dann aber geht Fall 3) und 4) in Fall 1) und 2) \"uber, denn die $\sigma,\tau$ entsprechenden Zahlen werden: $\sigma'=\sigma-i_{\tau-\sigma}+\tau-\sigma$, $\tau'=\tau-i_{\tau-\sigma}=\sigma'$; und bei Fortsetzung des Verfahrens wird $\tau'<\sigma'$. Es ist also im Schlu\ss ausdruck von (25.) das erste Summenzeichen zu ersetzen durch $\sum_{\alpha=\tau-\sigma}^{\tau,\rho}$.

Analog ergibt sich aus (21.) die dualistische Zerlegungsidentit\"at, bei der eine Reihe $q_\rho$ in die einzelnen Reihen $v^{(1)},\ldots,v^{(\rho)}$ zerlegt wird. Setzen wir
\begin{equation}
 \dot p_{\tau-\beta}^{(i_1\ldots i_\beta)}\sim v^{(i_1)}\cdots v^{(i_\beta)}q_{n-\tau};\qquad
 \dot q_{\rho-\beta}^{(i_1\ldots i_\beta)}=v^{(i_{\beta+1})}\cdots v^{(i_\rho)};\qquad
 q_\rho=v^{(1)}\cdots v^{(\rho)},
\tag{29}
\end{equation}
so kommt:
\begin{equation}
 \pair{q_\rho A^\sigma}{p_\tau B_{\rho+\sigma-\tau}}
 =\sum_{\substack{0,\tau-\sigma}}^{\substack{\tau,\rho}}
   \sum_{i_\beta=i_{\beta-1}+1}^{\rho}\eps
\pair{\dot q_{\rho-\beta}^{(i_1\ldots i_\beta)}B^{n-\rho-\sigma+\tau}}{A_{n-\sigma}\dot p_{\tau-\beta}^{(i_1\ldots i_\beta)}} ,
\tag{30}
\end{equation}
wo der Summationsindex $\beta$ vom gr\"o\ss eren der unteren Werte bis zum kleineren der oberen l\"auft. Alle Glieder einer Einzelsumme in (25.) und (30.) sind Polaren einer einzigen Form, entstanden durch die Polarprozesse
\[
\left(\frac{\partial}{\partial p_{\tau-\alpha}}\middle|p_{\tau-\alpha}^{(i_1\ldots i_\alpha)}\right),\qquad
\left(q_{\rho-\alpha}^{(i_1\ldots i_\alpha)}\middle|\frac{\partial}{\partial q_{\rho-\alpha}}\right).
\]
Um dies zum Ausdruck zu bringen, bezeichnen wir mit $\Pi$ ein Aggregat solcher Polaroperationen, multipliziert mit Konstanten, und erhalten dann an Stelle von (25.) und (30.) die Relation:
\begin{equation}
 \pair{q_\rho A^\sigma}{p_\tau B_{\rho+\sigma-\tau}}
 =\sum_{\substack{0,\tau-\sigma}}^{\substack{\tau,\rho}}
 \Pi\pair{q_{\rho-\alpha}B^{n-\rho-\sigma+\tau}}{A_{n-\sigma}p_{\tau-\alpha}}.
\tag{31}
\end{equation}

\section*{\S\ 5. Die Zerlegungsidentit\"aten in expliziter Form.}
Spezialisieren wir in (25.) im Falle 4) $\tau$ zu $\sigma+\rho$, so kommt:
\begin{equation}
 \pair{q_\rho A^\sigma}{y^{(1)}\cdots y^{(\rho+\sigma)}}
 =\sum_{i_\rho=i_{\rho-1}+1}^{\tau}
 \eps_{i_\rho}\pair{q_\rho}{y^{(i_1)}\cdots y^{(i_\rho)}}
 \pair{A^\sigma}{y^{(i_{\rho+1})}\cdots y^{(i_{\rho+\sigma})}},
\tag{32}
\end{equation}
wo $\eps_{i_\rho}=(-1)^{\delta_{i_\rho}}$, $\delta_{i_\rho}=\rho(\rho+1)/2+i_1+\cdots+i_\rho$ wird. Dies ist eine allgemeinere Form des Produktsatzes f\"ur Matrizen als (16.) und (17.) und ergibt sich durch Laplacesche Entwicklung der Determinante der Produktenmatrix
\[
 \sum\pm(v^{(1)}y^{(1)})\cdots(v^{(\rho)}y^{(\rho)})(a^{(1)}y^{(\rho+1)})\cdots(a^{(\sigma)}y^{(\rho+\sigma)}).
\]
Der allgemeinere Produktsatz ist also aus den Identit\"aten (20.), bzw. (21.) zusammengesetzt, und das Herausgreifen der speziellen Form (16.), (17.) ist dadurch erkl\"art.

Die Relation (32.) gibt nun das Mittel zur expliziten Darstellung der Zerlegungsidentit\"aten. Wir betrachten dazu die in (31.) auftretenden Formen als Grundformen; d. h. wir f\"uhren die Substitution (11.) aus, die an der expliziten Darstellung in den Reihen $A,B$ nichts \"andert, da nach (8.) die neu eingef\"uhrten Reihen $R$ sich durch die Reihen $A,B$ selbst und nicht durch deren Teilreihen $a^{(1)}a^{(2)}\ldots$, $b^{(1)}b^{(2)}\ldots$ ausdr\"ucken lassen. Sei also
\begin{equation}
 \pair{q_{\rho-\alpha}B^{n-\rho-\sigma+\tau}}{A_{n-\sigma}p_{\tau-\alpha}}
 =\bigl(R_{\rho-\alpha}p_{n-\rho+\alpha}\bigr)
  \pair{R^{\tau-\alpha}}{p_{\tau-\alpha}},
\tag{33}
\end{equation}
so kommt f\"ur die dem Index $\alpha$ entsprechende Einzelsumme von (25.):
\begin{equation}
\begin{aligned}
&\sum \eps_{i_\alpha}\bigl(R_{\rho-\alpha}p_{n-\rho}y^{(i_1)}\cdots y^{(i_\alpha)}\bigr)
 \pair{R^{\tau-\alpha}}{y^{(i_{\alpha+1})}\cdots y^{(i_\tau)}}\\
&\quad=\sum \eps_{i_\alpha}\bigl([R_{\rho-\alpha}p_{n-\rho}]^\alpha y^{(i_1)}\cdots y^{(i_\alpha)}\bigr)
 \pair{R^{\tau-\alpha}}{y^{(i_{\alpha+1})}\cdots y^{(i_\tau)}}\\
&\quad=\eps\pair{[R_{\rho-\alpha}p_{n-\rho}]^\alpha R^{\tau-\alpha}}{p_\tau}
 =\eps\pair{R^{\tau-\alpha}q_{n-\tau}}{R_{\rho-\alpha}p_{n-\rho}}.
\end{aligned}
\tag{34}
\end{equation}
Damit ist nun in der Tat ein expliziter Schlu\ss ausdruck gegeben, und die symmetrische Form, in der $q_\rho$ und $p_\tau$ eingehen, zeigt, da\ss\ (30.) zum gleichen Schlu\ss ausdruck f\"uhrt, der also unabh\"angig ist von der urspr\"unglichen Zerlegung. Ein Unterschied zwischen (25.) und (30.) besteht nur, solange die Reihen $y$ und $v$ selbst\"andige Bedeutung haben, die Zerlegungsidentit\"at also nach unserer Definition in eine zerlegbare Identit\"at \"ubergeht. Um zu Identit\"aten mit mehr Symbolreihen zu gelangen, hat man nach \S\ 3 (Schlu\ss) nur eine Reihe $q_\rho$ in $q_{\rho_1}C^{\sigma_1}$, bzw. $p_\tau$ in $p_{\tau_1}C_{\tau-\tau_1}$ aufzul\"osen. Die dann entstehenden Ausdr\"ucke:
\begin{equation}
 \pair{q_{\rho_1}C^{\sigma_1}}{[R^{\tau-\alpha}q_{n-\tau}]_\lambda R_{\rho_1+\sigma_1-\alpha-\lambda}},
 \qquad
 \pair{[R_{\rho-\alpha}p_{n-\rho}]^\lambda R^{\tau-\alpha}}{p_{\tau_1}C_{\tau-\tau_1}},
\tag{35}
\end{equation}
sind nun genau vom Typus der bei zwei Symbolreihen auftretenden, lassen also bei einer (33.) entsprechenden Substitution wieder explizite Darstellung zu, so da\ss\ sich durch Wiederholung explizite Darstellung bei beliebig vielen Reihen ergibt. Bemerkt mag noch werden, da\ss\ je die erste und letzte Summe in (25.) bzw. (30.) die explizite Darstellung ohne Anwendung von (33.) allein durch die Formel (32.) ergibt, oder sich \"uberhaupt auf ein einziges, alle Reihen explizit enthaltendes Glied beschr\"ankt. Die durch Aufl\"osung von $q_\rho$ in $q_{\rho_1}C^{\sigma_1}$ usf. aus (32.) entstehende Relation kann als der Produktsatz in seiner allgemeinsten Form aufgefa\ss t werden.

Zusammenfassend werden wir zeigen:

\emph{Satz III:} Mit den Identit\"aten
\begin{equation}
 \pair{q_\rho A^\sigma}{p_\tau B_{\rho+\sigma-\tau}}
 =\sum_{\substack{0,\tau-\sigma}}^{\substack{\tau,\rho}}
 \eps\pair{R^{\tau-\alpha}q_{n-\tau}}{R_{\rho-\alpha}p_{n-\rho}},
\tag{36}
\end{equation}
wo die Reihen $R$ durch (33.) bestimmt sind, und mit den durch Aufl\"osung von $q_\rho,p_\tau$ in $q_{\rho_1}C^{\sigma_1}$ usf. daraus entstehenden ist die Gesamtheit der unzerlegbaren Identit\"aten ersch\"opft.

In der Tat sind die entstehenden Identit\"aten die allgemeinsten ihrer Art, da die einzelnen Reihen in bezug auf Gewicht und Anzahl keiner Beschr\"ankung mehr unterliegen, wie dies noch bei (20.), (21.) der Fall war. Es existieren aber auch keine weiteren Identit\"aten zwischen diesen Reihen; denn bei Aufl\"osung einer Reihe $p_\tau$ in $y^{(1)}\cdots y^{(\tau)}$ sind die allgemeinsten Identit\"aten aus (20.), also nach \S\ 3 (Schlu\ss) aus (20a.) zusammengesetzt; und zwar ist die im Anschlu\ss\ an (16.) diskutierte Zusammensetzung genau die in \S\ 4 durchgef\"uhrte. Den \"Ubergang zu beliebig vielen Reihen liefert aber, wie die Diskussion von (20a.) zeigt, die Anwendung immer derselben Operation auf die durch Aufl\"osung von $q_\rho$ in $q_{\rho_1}C^{\sigma_1}$ usf. entstehenden Ausdr\"ucke. Es folgt daraus auch, da\ss\ der direkte \"Ubergang von (20.) an Stelle von (20a.) nichts Neues liefert; dies l\"a\ss t sich rechnerisch leicht best\"atigen, indem man einmal eine Reihe $q_\rho$ in (25.) spezialisiert, das andere Mal nach der Methode von \S\ 4 den \"Ubergang von (20.) macht; es entstehen beide Male die gleichen Summen, die durch Anwendung des allgemeinen Produktsatzes (siehe oben) in die aus (36.) sich ergebenden Identit\"aten \"ubergehen. Die Identit\"aten (20.), (21.) entsprechen den Spezialf\"allen $\tau=1$, $\rho=1$; es treten dann nur je zwei Einzelsummen auf, also einer obigen Bemerkung entsprechend explizite Darstellung ohne die Substitution (33.), was mit dem fr\"uheren \"ubereinstimmt.

Zum Schlu\ss\ seien noch zwei Spezialf\"alle der Zerlegungsidentit\"at kurz erw\"ahnt. Setzen wir in (31.): $\tau=\sigma$, $\rho=n-\sigma$ und bezeichnen die Reihe $p_\tau$ als $p'_\sigma\sim q'_{n-\sigma}$, so kommt:
\begin{equation}
\begin{aligned}
\pair{A^\sigma}{p_\sigma}\pair{B^\sigma}{p'_\sigma}
-\pair{A^\sigma}{p'_\sigma}\pair{B^\sigma}{p_\sigma}
=\sum_{\alpha=1}^{\sigma,n-\sigma}
\Pi\pair{q_{n-\sigma-\alpha}B^\sigma}{A_{n-\sigma}p_{\sigma-\alpha}}.
\end{aligned}
\tag{37}
\end{equation}
Das ist die Clebschsche Formel\footnote{Clebsch, a. a. O., S. 25--32. Clebsch weist nach, da\ss\ die einzelnen Glieder der rechten Seiten nur von den Reihen $A,B$ abh\"angen; die explizite Darstellung w\"urde sich durch Anwendung von (32.) auf seine Ausdr\"ucke $\Phi_h$ ergeben.} zur Entwicklung einer Form mit zwei kogredienten Reihen $p_\sigma,p'_\sigma$ nach Elementarkovarianten. Die zu einer Form $\pair{A^\sigma}{p_\sigma}^m\pair{B^\sigma}{p'_\sigma}^n$ geh\"orenden Elementarkovarianten sind dann die folgenden:
\begin{equation}
 \left\{\sum_{\alpha=1}^{\sigma,n-\sigma}
 \pair{q_{n-\sigma-\alpha}B^\sigma}{A_{n-\sigma}p_{\sigma-\alpha}}
 \right\}^{\rho}
 \pair{A^\sigma}{p_\sigma}^{m-\rho}\pair{B^\sigma}{p'_\sigma}^{n-\rho},
 \qquad (\rho=0,1,\ldots,m,n).
\tag{38}
\end{equation}
Sie gehen, wie wir kurz sagen wollen, aus der Grundform durch ,,kogrediente Faltung vom Defekt $\rho$`` (vgl. \S\ 2) hervor.

Setzen wir zweitens: $\tau=\sigma-\lambda$, $\rho=n-\sigma-\lambda$; $B^\sigma=A^\sigma$, so kommt:
\begin{equation}
\begin{aligned}
\pair{q_{n-\sigma-\lambda}A^\sigma}{A_{n-\sigma}p_{\sigma-\lambda}}
&=(-1)^\lambda\pair{q_{n-\sigma-\lambda}A^\sigma}{A_{n-\sigma}p_{\sigma-\lambda}}\\
&\quad+(-1)^{(n-\sigma)(\sigma-\lambda)}
 \sum_{\alpha=1}^{\sigma-\lambda,n-\sigma-\lambda}
 \Pi\pair{q_{n-\sigma-\lambda-\alpha}A^\sigma}{A_{n-\sigma}p_{\sigma-\lambda}}.
\end{aligned}
\tag{39}
\end{equation}
Es sind also, wie in \S\ 2 bemerkt, die die Symbolreihen charakterisierenden Bedingungen (12.) f\"ur ungeraden Defekt $\lambda$ eine Folge der Bedingungen von h\"oherem Defekt. F\"ur eine Linearform $\pair{A^\sigma}{p_\sigma}=\pair{B^\sigma}{p_\sigma}$ ergibt sich aus (39.), da\ss\ sich ihre irreduzibeln invarianten Bildungen, die in den Koeffizienten quadratisch sind, auf solche von geradem Defekt reduzieren; das stimmt mit der bekannten Tatsache \"uberein, da\ss\ eine Linearform im bin\"aren und tern\"aren Gebiet keine invarianten Bildungen besitzt.

Bemerkt sei, da\ss\ die allgemeinen Identit\"aten urspr\"unglich abgeleitet sind unter der Annahme, da\ss\ die einzelnen Reihen den Bedingungen (12.) gen\"ugen. Da indes diese einzelnen Reihen nur linear in die Identit\"aten eingehen, so gelten letztere auch f\"ur die allgemeineren, den Bedingungen (12.) nicht gen\"ugenden Reihen.

\section*{\S\ 6. Explizite Darstellung der invarianten Bildungen.\\Faltungs- und Differentiationsprozesse.}
Die Resultate der vorigen Paragraphen erlauben uns, den in \S\ 2 ausgesprochenen Satz von der symbolischen Darstellbarkeit (1. Fundamentalsatz) durch Zuf\"ugung der expliziten Darstellbarkeit zum Abschlu\ss\ zu bringen, indem wir zeigen:

\emph{Satz IV:} ,,Alle invarianten Bildungen beliebiger Grundformen lassen sich explizit durch Matrizenprodukte darstellen; d. h. in die Schlu\ss{}ausdr\"ucke gehen au\ss er Variabelnreihen $p_\rho$ nur die Reihen $S^\sigma,\ldots,T^\tau$ der urspr\"unglichen Formen oder die durch Substitution (9.) oder (11.) aus ihnen abgeleiteten Reihen $P,Q,R$ ein.``

Zum Beweise nehmen wir an, eine invariante Bildung m\"oge au\ss er von Variabelnreihen von den Reihen $S^\sigma,\ldots,T^\tau$ abh\"angen; und zwar seien diese Reihen gen\"ugend weit in einzelne Reihen: $S^\sigma=S^{\sigma_1}\cdots S^{\sigma_i},\ldots,T^\tau=T^{\tau_1}\cdots T^{\tau_k}$ aufgel\"ost, um eine Darstellung durch Determinanten zu erm\"oglichen. Die Reihen seien in einer an sich willk\"urlichen, aber bestimmt festgelegten Reihenfolge $S^\sigma,\ldots,T^\tau$ geordnet. Wir greifen nun ein solches Aggregat von Determinanten heraus, das von der ersten Gesamtreihe $S^\sigma=S^{\sigma_1}\cdots S^{\sigma_i}$, nicht nur von einzelnen der Teilreihen, abh\"angt. Diese Abh\"angigkeit ist aber nach \S\ 3--\S\ 5 eine Folge der allgemeinen Identit\"aten. L\"osen wir nun eine Variabelnreihe $p_\rho$ in die Reihen $y,v$, oder eine der sp\"ateren Reihen $T$ in die einzelnen Reihen $t$, bzw. $\tau$ auf, so sind die allgemeinsten Identit\"aten aus (20.), (21.) zusammengesetzt. Diese letzteren Identit\"aten f\"uhren aber immer zu einem die Reihe $S^\sigma=S^{\sigma_1}\cdots S^{\sigma_i}$ explizit enthaltenden Ausdruck, der linken Seite von (20.), (21.); man \"uberzeugt sich n\"amlich durch (19.), da\ss\ erst die vollst\"andige rechts stehende Summe -- nicht aber ein Teil derselben, der einem Gliede von (20a.) entspricht -- die Eigenschaft besitzt, nur von $S^\sigma$ abzuh\"angen. Bei Fortsetzung des Verfahrens bleiben alle einmal explizit eingehenden Reihen auch explizit erhalten; die Anwendung der Identit\"aten kann nur m\"oglicherweise ein Zusammenfassen mehrerer Reihen zu einer einzigen, d. h. die Substitution (9.), verlangen. Ist schlie\ss lich nur noch eine einzige Reihe $T^\tau$ auf die explizite Form zu bringen, so sind zwei F\"alle m\"oglich. Es kann noch eine freie, d. h. nicht durch Substitution (9.) mit andern Reihen verbundene, Variabelnreihe auftreten; durch deren Zerlegung in Teilreihen $y$ oder $v$ l\"a\ss t sich das Verfahren zum Abschlu\ss\ bringen; es entstehen, wie in (31.), Polaren einer oder mehrerer, alle Reihen explizit enthaltenden Formen. Tritt keine freie Variabelnreihe mehr auf, so erkennen wir nach der Art, wie sie entstanden sind, in der Gesamtheit der Ausdr\"ucke, die die einzelnen Reihen $t$ oder $\tau$ enthalten, aber von der Reihe $T^\tau$ abh\"angen, die Glieder einer Zerlegungsidentit\"at; diese aber lassen nach \S\ 5 unter Anwendung der Substitution (11.) stets die explizite Darstellung in allen Reihen zu. Damit ist der obige Satz allgemein bewiesen. Bemerkt mag noch werden, da\ss\ dann, wenn nur zwei Symbolreihen auftreten, die Substitution (11.) nie eintritt; es k\"onnen n\"amlich -- wegen $\sigma,n-\sigma,\tau,n-\tau<n$ -- nur dann keine Variabelnreihen eingehen, wenn die beiden Symbolreihen in einer einzigen Determinante vereinigt, also explizit auftreten; sobald aber Variabelnreihen eingehen, zeigen (34.) und (33.), da\ss\ die Substitution (11.) \"uberfl\"ussig wird.

Aus dem eben Gesagten folgt, da\ss\ es zur Erzeugung aller invarianten Bildungen gen\"ugt, die Symbolreihen unter Zuf\"ugung von Variabelnreihen zu allen m\"oglichen Matrizenprodukten zu vereinigen. Das allgemeinste Matrizenprodukt ist nun von der Form:
\begin{equation}
 \pair{P^{\mu_1}\cdots P^{\mu_i}}{Q_{\nu_1}\cdots Q_{\nu_k}},\qquad \sum\mu=\sum\nu,
\tag{40}
\end{equation}
wo die $P$ und $Q$ Reihen der urspr\"unglichen Formen sein k\"onnen, oder durch eine Reihenfolge von Substitutionen (9.) oder (11.) aus den urspr\"unglichen Reihen hervorgegangen, oder endlich Variabelnreihen. Die Ausdr\"ucke (40.) aber lassen sich erzeugen durch sukzessive Vereinigung je zweier Symbolreihen zu einem Matrizenprodukt, unter Zuhilfenahme von Variabelnreihen; denn man erh\"alt aus
\[
 \pair{q_{n-\mu-\lambda}P^\mu}{Q_{\nu_1}p_{n-\nu_1-\lambda}}
 =\bigl(R_\lambda Q_{\nu_1}p_{n-\nu_1-\lambda}\bigr)
\]
durch Vereinigung mit der Reihe $Q_{\nu_2}$ einen Ausdruck
\[
 \pair{[R_\lambda Q_{\nu_1}]^{n-\lambda-\nu_1}}{Q_{\nu_2}p_{n-\nu_1-\nu_2-\lambda}}
 =\pair{q_{n-\mu-\lambda}P^\mu}{Q_{\nu_1}Q_{\nu_2}p_{n-\nu_1-\nu_2-\lambda}},
\]
also durch Fortsetzung des Verfahrens einen Ausdruck (40.). Sind dabei $P$ und $Q$ nicht die Reihen der urspr\"unglichen Formen, so zeigt (9.) und (11.) in Verbindung mit dem eben Gezeigten, da\ss\ sie gleichfalls durch eine Reihenfolge von Vereinigungen je zweier Symbolreihen entstanden sind. Es folgt daraus:

\emph{Satz V:} Der Proze\ss\ zur Erzeugung aller invarianten Bildungen, der Faltungsproze\ss, besteht in der sukzessiven Vereinigung je zweier Symbolreihen zu einem Matrizenprodukt, unter Zuhilfenahme von Variabelnreihen.

Aus zwei Symbolreihen $S^\sigma,T^\tau$, wo $\sigma\geqq\tau$, lassen sich nun die folgenden Matrizenprodukte bilden:
\begin{align}
&\pair{q_{n-\sigma-\lambda}S^\sigma}{T_{n-\tau}p_{\tau-\lambda}}, &&(\lambda=1,2,\ldots,\tau,n-\sigma),\tag{41}\\
&\pair{q_{n-\tau-\mu}T^\tau}{S_{n-\sigma}p_{\sigma-\mu}}, &&(\mu=\sigma-\tau+\lambda),\tag{42}\\
&\pair{q_{n-\tau-\nu}T^\tau}{S_{n-\sigma}p_{\sigma-\nu}}, &&(\nu=1,2,\ldots,\sigma-\tau).
\tag{43}
\end{align}
Aus (31.) ergeben sich aber f\"ur (42.) und (43.) die Werte:
\begin{align}
&\pair{q_{n-\sigma-\lambda}S^\sigma}{T_{n-\tau}p_{\tau-\lambda}}
 +\sum_{\alpha}\Pi\pair{q_{n-\sigma-\lambda-\alpha}S^\sigma}{T_{n-\tau}p_{\tau-\lambda-\alpha}},\tag{42a}\\
&\Pi(q_{n-\sigma}S^\sigma)(T_{n-\tau}p_\tau)
 +\sum_{\beta}\Pi\pair{q_{n-\sigma-\beta}S^\sigma}{T_{n-\tau}p_{\tau-\beta}}.
\tag{43a}
\end{align}
Es lassen sich also die Formen (42.) linear ausdr\"ucken durch (41.) und Polaren von (41.), die Formen (43.) durch Polaren der Grundform und der Formen (41.). Ebenso ergibt sich aus (31.) die lineare Abh\"angigkeit der Formen (41.) von (42.) und Polaren von (42.). Danach sind aber nach der Clebschschen Definition\footnote{Clebsch, a. a. O. \S\ 7.} die Formen (42.) den Formen (41.) \"aquivalent, w\"ahrend (43.) auf die Grundform zur\"uckf\"uhrt. Der erzeugende Proze\ss\ ist also gleichberechtigt durch (41.) oder (42.) gegeben; wir wollen den in bezug auf die Zahl $\lambda$ einfacheren Proze\ss\ (41.) als Faltungsproze\ss\ definieren. Die Zahl $\lambda$, der Defekt der Faltung (vgl. \S\ 2), gibt die Anzahl der zum Determinantenprodukt fehlenden Reihen an, und zwar in demjenigen Matrizenprodukt, das bei gegebener Faltung die Maximalzahl von Reihen enth\"alt. Da $\lambda$ nur bis zum kleineren der beiden Werte $\tau,n-\sigma$ l\"auft, wo $\sigma\geqq\tau$, so ist:
\begin{equation}
 \lambda\leqq\frac n2,\ n\text{ gerade};\qquad
 \lambda\leqq\frac{n-1}{2},\ n\text{ ungerade}.
\tag{44}
\end{equation}
Die Zur\"uckf\"uhrung von (42.) und (43.) auf (41.) gilt auch dann noch, wenn durch weitere Faltung die Reihen $p,q$ in Symbolreihen \"ubergegangen sind; denn nach (36.) kommt:
\[
 \pair{A^{n-\tau-\varkappa}T^\tau}{S_{n-\sigma}B_{\sigma-\varkappa}}
 =\sum_{\alpha=0,\varkappa-\sigma+\tau}^{\tau,n-\sigma}
 \eps\pair{R^{\tau-\alpha}B^{n-\sigma+\varkappa}}{A_{\tau+\varkappa}R_{n-\sigma-\alpha}},
\]
wo die Reihen $R$ nach (33.) aus $S$ und $T$ hervorgegangen sind. Man erh\"alt also diese nur Symbolreihen enthaltenden Formen durch Faltung von $\pair{q_\tau A^{n-\tau-\varkappa}}{B_{\sigma-\varkappa}p_{n-\sigma}}$ bzw. $\pair{q_{\tau-\alpha}B^{n-\sigma+\varkappa}}{A_{\tau+\varkappa}p_{n-\sigma-\alpha}}$ -- je nachdem $\sigma-\tau\gtreqless 2\varkappa$ -- mit der Grundform und den Formen (41.).

Aus dem symbolischen Faltungsproze\ss\ gewinnt man in bekannter Weise die invarianten Differentiationsprozesse, wenn man nur die Symbolreihen $S^\sigma,T_{n-\tau}$ durch die Reihen von Differentiationssymbolen $\frac{\partial}{\partial p_\sigma}$, $\frac{\partial}{\partial q_{n-\tau}}$ ersetzt, wobei bekanntlich jedem Produkte $\frac{\partial}{\partial p_{i_1\ldots i_\sigma}}\frac{\partial}{\partial q_{j_1\ldots j_{n-\tau}}}$ die reale Bedeutung $\frac{\partial^2}{\partial p_{i_1\ldots i_\sigma}\partial q_{j_1\ldots j_{n-\tau}}}$ zukommt. Da die Reihen $\frac{\partial}{\partial p_\sigma}$, $\frac{\partial}{\partial q_{n-\tau}}$ den Reihen $S^\sigma,T_{n-\tau}$ kogredient sind, gen\"ugen sie denselben Identit\"aten wie diese, also speziell auch den zwischen (42.) und (42a.), (43.) und (43a.) bestehenden. Zusammenfassend erhalten wir:

\emph{Satz VI:} Alle invarianten Bildungen beliebiger Grundformen entstehen aus diesen durch wiederholte Anwendung des symbolischen Faltungsprozesses (41.) oder auch der invarianten Differentiationsprozesse:
\begin{equation}
 \left(q_{n-\sigma-\lambda}\frac{\partial}{\partial p_\sigma}\middle|\frac{\partial}{\partial q_{n-\tau}}p_{\tau-\lambda}\right),
 \qquad(\sigma\geqq\tau,\ \lambda=1,2,\ldots,\sigma,n-\tau).
\tag{45}
\end{equation}
Es ist noch zu bemerken, da\ss\ bei jeder Faltung (41.) das Gesamtgewicht der Form, d. h. die Summe der Gewichte der einzelnen Variabelnreihen (vgl. \S\ 1), abnimmt um die positive Zahl $2(\lambda^2+\lambda(\sigma-\tau))$. Daraus folgt, unabh\"angig vom allgemeinen Endlichkeitsbeweis, die Endlichkeit der Formenreihe, d. h. der Gesamtheit der in den Koeffizienten linearen invarianten Bildungen einer gegebenen Grundform.\footnote{Vgl. Clebsch, a. a. O. \S\ 11 u. 12. Endlichkeit des Systems von Elementarkovarianten.}

Bemerkt sei weiter, da\ss\ Satz VI auch noch gilt f\"ur Formen, die den Bedingungen (12.) nicht gen\"ugen. Es l\"a\ss t sich n\"amlich jeder Faktor von (3.): $\pair{S^\rho}{p_\rho}^{m_\rho}$ ersetzen durch ein Produkt von $m_\rho$ Linearfaktoren $\pair{A^\rho}{p_\rho}\pair{B^\rho}{p_\rho}\cdots\pair{M^\rho}{p_\rho}$; wodurch die invarianten Bildungen \"ubergehen in solche eines Systems von Linearformen, die erst nachtr\"aglich einander gleich gesetzt zu werden brauchen. F\"ur jede einzelne Linearform aber sind die Bedingungen (12.), die bei Einf\"uhrung von Differentialsymbolen nur Differentialquotienten 2. Ordnung enthalten, stets erf\"ullt.

